\chapter{Complexity — Meyer Machine Run}
%%% generated by the blueprint pipeline from TCSlib.Complexity.CircuitComplexity.MeyerMachineRun
\section{Overview}

Declarations of the Lean module \texttt{TCSlib.Complexity.CircuitComplexity.MeyerMachineRun}.

\section{Declarations}

\begin{definition}[Configuration of the Simulating Machine]
\lean{Complexity.Meyer.lcfg}
\label{Complexity.Meyer.lcfg}
\leanok
\uses{Complexity.Meyer.St, Complexity.Meyer.tabTM, Complexity.Meyer.tcfg, Turing.Cfg, Turing.FinTM.bufferTape}
Let $M$ be a finite Turing machine over the Boolean alphabet, let $x$ and $y$ be Boolean
strings, and let $c$ be a configuration of $M$ on input $x$. Given a state $s$ of the
simulating machine, an input-head position $P$ on input $y$, two auxiliary tapes
$\tau_C,\tau_O:\bbr\to\text{Option}(\text{Bool})$ with head positions
$z_C,z_O\in\bbr$, define a configuration on input $y$ with state $s$,
input-head position $P$, the work tapes and head positions of $c$ followed by the tape
storing $x$ and the tapes $\tau_C,\tau_O$ with head positions
$\text{inputPos}(c)-1,z_C,z_O$, respectively, and an empty output tape.
\end{definition}

\begin{theorem}[Equality of two simulated configurations]
\lean{Complexity.Meyer.acfg_eq_lcfg}
\label{Complexity.Meyer.acfg_eq_lcfg}
\leanok
\uses{Complexity.Meyer.St, Complexity.Meyer.acfg, Complexity.Meyer.lcfg, Turing.FinTM.bufferTape, Turing.MultiTapeTM.initCfg}
\difficulty{3}
Let $M$ be a finite-state Turing machine with $k$ work tapes, let $y$ and $x$ be Boolean
words, let $st$ be a state of the simulator built from $M$, and let $P$ be a valid
input-head position for $y$. A simulator configuration with state $st$, input-head
position $P$, three auxiliary tapes storing $x$, $t$, and $a$ from cell $0$ onward, all
auxiliary heads at cell $0$, all simulated work tapes blank with their heads at the
origin, and empty output equals the lifted configuration obtained from $M$'s initial
configuration on $x$, with the same simulator state and input-head position, the
auxiliary tapes storing $x$, $t$, and $a$, and the auxiliary heads at cell $0$. Here $t$
and $a$ are arbitrary Boolean words.
\end{theorem}

\begin{theorem}[One-step simulation configuration identity]
\lean{Complexity.Meyer.lcfg_step}
\label{Complexity.Meyer.lcfg_step}
\leanok
\uses{Complexity.Meyer.St, Complexity.Meyer.lcfg, Complexity.Meyer.tabTM, Complexity.Meyer.tabTM_step, Complexity.Meyer.tabTr, Turing.Action.apply, Turing.Cfg, Turing.Cfg.inputSymbol, Turing.Cfg.workTapeSymbols, Turing.FinTM.bufferTape_inputSymbol, Turing.MultiTapeTM.step}
\difficulty{4}
Let $M$ be a finite Turing machine over the binary alphabet, let $x,y$ be finite binary
strings, and choose a simulator state $st$, an input-head position $P$ for $y$, a
configuration $c$ of $M$ on $x$, auxiliary tape contents $\tau_C,\tau_O$, and auxiliary
head positions $z_C,z_O$. The simulator machine's one-step successor of the lifted
configuration is exactly the result of applying its transition action at that
configuration, using the scanned input symbol and the tuple consisting of the symbols
under the original work-tape heads and the contents of the auxiliary tapes at their
respective head positions.
\end{theorem}

\begin{theorem}[Effect of a lifted action]
\lean{Complexity.Meyer.lcfg_apply_aux}
\label{Complexity.Meyer.lcfg_apply_aux}
\leanok
\uses{Complexity.Meyer.St, Complexity.Meyer.idle, Complexity.Meyer.lcfg, Complexity.Meyer.tact, Complexity.Meyer.tact_apply, Complexity.Meyer.tcfg, Complexity.Meyer.wr, Turing.Action.apply, Turing.Cfg, Turing.FinTM.bufferTape, Turing.moveInputPos_zero}
\difficulty{5}
Let $M$ be a finite Turing machine over the binary alphabet, let $x,y$ be finite binary
strings, and choose a simulator state, a simulator input-head position $P$, a
configuration $c$ of $M$ on input $x$, two work-tape contents $\tau_C,\tau_O$ with head
positions $z_C,z_O$, an action $a_w$ on the simulator's three auxiliary tapes, an
optional output bit $o$, and an optional successor simulator state $q$. If the first
auxiliary tape is left unchanged with no movement, applying the simulator action with no
input-head movement produces the simulator configuration with state $q$ and position
$P$, the original work tapes and heads of $c$, auxiliary contents given by writing the
prescribed values of $a_w$ at $z_C,z_O$, auxiliary head positions $z_C$ and $z_O$
shifted by their prescribed movements, and output consisting of $o$ when present and the
empty string otherwise.
\end{theorem}

\begin{theorem}[Original Work-Tape Readings]
\lean{Complexity.Meyer.mreads}
\label{Complexity.Meyer.mreads}
\leanok
\uses{Turing.Cfg.workTapeSymbols}
\difficulty{1}
Let $M$ be a Turing machine with $k$ work tapes and Boolean tape symbols, and let $x$ be
a Boolean input with configuration $c$. For any function $g$ assigning an optional
Boolean symbol to each of three additional tapes, the symbols currently read on the
first $k$ tapes of the combined $k+3$ tape collection are exactly the symbols read by
the work tapes of $c$.
\end{theorem}

\begin{theorem}[Entering the Write-Decode State]
\lean{Complexity.Meyer.dec_under}
\label{Complexity.Meyer.dec_under}
\leanok
\uses{Complexity.Meyer.aux_read1, Complexity.Meyer.lcfg, Complexity.Meyer.lcfg_apply_aux, Complexity.Meyer.lcfg_step, Complexity.Meyer.tabTM, Complexity.Meyer.tabTr, Complexity.TimeHierarchy.OutReg, Turing.MultiTapeTM.step}
\difficulty{3}
Let $M$ be a finite Boolean Turing machine, let $y,x$ be Boolean words, let $bits$ be a
selector-bit assignment for $M$, let $P$ be a valid input-head position for $y$, let $c$
be a configuration of $M$ on input $x$, and let $\tau_C,\tau_O$ be integer-indexed
Boolean buffers with head positions $z_C,z_O$. If the $C$ buffer is blank at $z_C$, then
one step of the simulator from its decode state yields its write-decode state, retaining
$P$, $c$, both buffers, and both head positions.
\end{theorem}

\begin{theorem}[Zero-cell counter update]
\lean{Complexity.Meyer.dec_zero}
\label{Complexity.Meyer.dec_zero}
\leanok
\uses{Complexity.Meyer.aux_read1, Complexity.Meyer.lcfg, Complexity.Meyer.lcfg_apply_aux, Complexity.Meyer.lcfg_step, Complexity.Meyer.tabTM, Complexity.Meyer.tabTr, Complexity.TimeHierarchy.OutReg, Turing.MultiTapeTM.step}
\difficulty{3}
Let $M$ be a finite-state Turing machine with Boolean tape symbols, and let $y,x$ be
Boolean words. Choose a bit vector encoding a transition-table selector, a simulated
input-head position, and a configuration $c$ of $M$ on $x$. Let the counter tape have
contents $\tau_C:\mathbb{Z}\to\mathrm{Option}(\mathrm{Bool})$ and head position
$z_C\in\mathbb{Z}$, and let $\tau_O$ and $z_O$ specify the other auxiliary tape. If the
counter tape contains the Boolean symbol false at $z_C$, then one step of the tabulating
machine changes the simulated configuration to the same decoding state, writes true at
$z_C$, and moves the counter head to $z_C+1$, leaving the other specified data
unchanged.
\end{theorem}

\begin{theorem}[Return Transition at a Nonblank Cell]
\lean{Complexity.Meyer.ret_some}
\label{Complexity.Meyer.ret_some}
\leanok
\uses{Complexity.Meyer.aux_read1, Complexity.Meyer.lcfg, Complexity.Meyer.lcfg_apply_aux, Complexity.Meyer.lcfg_step, Complexity.Meyer.tabTM, Complexity.Meyer.tabTr, Complexity.TimeHierarchy.OutReg, Turing.MultiTapeTM.step}
\difficulty{3}
Let $M$ be a finite Turing machine over the binary alphabet, let $x$ and $y$ be bit
strings, and let $c$ be a configuration of $M$ on input $x$. Choose a selector bit
vector, an input-head position for the simulated machine on $y$, the simulated work
tapes and head positions, and two auxiliary tapes with head positions. If the first
auxiliary tape contains the bit $v$ at its head position $z_C$, then one step of the
tabulating machine from the corresponding return-state configuration leaves the state
and all specified data unchanged except that this tape's head moves from $z_C$ to
$z_C-1$.
\end{theorem}

\begin{theorem}[Return Transition at a Blank Cell]
\lean{Complexity.Meyer.ret_none}
\label{Complexity.Meyer.ret_none}
\leanok
\uses{Complexity.Meyer.aux_read1, Complexity.Meyer.lcfg, Complexity.Meyer.lcfg_apply_aux, Complexity.Meyer.lcfg_step, Complexity.Meyer.tabTM, Complexity.Meyer.tabTr, Complexity.TimeHierarchy.OutReg, Turing.MultiTapeTM.step}
\difficulty{3}
Let $M$ be a finite Turing machine over the binary alphabet, let $x$ and $y$ be bit
strings, and let $c$ be a configuration of $M$ on input $x$. Choose a selector bit
vector, an input-head position for the simulated machine on $y$, the simulated work
tapes and head positions, and two auxiliary tapes with head positions. If the first
auxiliary tape is blank at its head position $z_C$, then one step of the tabulating
machine from the corresponding return-state configuration is the corresponding
decision-state configuration, with that tape's head advanced from $z_C$ to $z_C+1$ and
the other specified data unchanged.
\end{theorem}

\begin{definition}[Next State After a Simulated Step]
\lean{Complexity.Meyer.fireNext}
\label{Complexity.Meyer.fireNext}
\leanok
\uses{Complexity.Meyer.St, Complexity.TimeHierarchy.OutReg}
For a finite Turing machine $M$ over the Boolean alphabet, a bit vector encoding a
selector, an input-head position, a configuration of $M$, two auxiliary tapes and their
head positions, and an output-register value $r$, the next simulating state determined
by a proposed successor state $q'$ and boundary tag $b$ is the return state carrying the
bit vector, $q'$, $r$, and $b$ when $q'$ exists; if $q'$ is absent, it is the
waiting-decision state carrying the bit vector, no state, $r$, and $b$.
\end{definition}

\begin{theorem}[One Simulated Machine Step During a Sweep]
\lean{Complexity.Meyer.dec_fire}
\label{Complexity.Meyer.dec_fire}
\leanok
\uses{Complexity.Meyer.aux_read0, Complexity.Meyer.aux_read1, Complexity.Meyer.fireNext, Complexity.Meyer.lcfg, Complexity.Meyer.lcfg_step, Complexity.Meyer.mreads, Complexity.Meyer.tabTM, Complexity.Meyer.tabTr, Complexity.Meyer.tact_apply, Complexity.Meyer.two, Complexity.TimeHierarchy.OutReg.ofList, Complexity.TimeHierarchy.OutReg.ofList_push, Complexity.TimeHierarchy.OutReg.push, Turing.Action.apply, Turing.Cfg.inputSymbol, Turing.Cfg.workTapeSymbols, Turing.FinTM.VirtualTag, Turing.FinTM.virtualMove, Turing.FinTM.virtualMove_correct, Turing.FinTM.virtualNextTag, Turing.MultiTapeTM.step, Turing.moveInputPos_zero}
\difficulty{6}
Let $M$ be a finite Boolean Turing machine, let $y,x$ be Boolean words, let $bits$ be a
selector-bit assignment for $M$, let $P$ be a valid input-head position for $y$, let $c$
be a configuration of $M$ on input $x$, and let
$\tau_C,\tau_O:\mathbb{Z}\to\{\mathrm{none},\mathrm{some} (\mathrm{false}),\mathrm{some}
(\mathrm{true})\}$ with head positions $z_C,z_O\in\mathbb{Z}$. If $c$ is in state $q$,
the Boolean tag $b$ correctly marks the input-head boundary of $c$, and the $C$ buffer
contains true at position $z_C$, then there is a tag $b'$ correctly marking the
input-head boundary after one step of $M$ such that one step of the simulator, started
in its decode state with the output register obtained from $c$'s output, produces the
lcfg configuration for that successor machine configuration. In that resulting
configuration, the $C$ buffer is updated at $z_C$ to false and its head moves to
$z_C-1$; the $O$ buffer and its head position are retained.
\end{theorem}

\begin{theorem}[Stored word at an occupied cell]
\lean{Complexity.Meyer.bufferTape_at}
\label{Complexity.Meyer.bufferTape_at}
\leanok
\uses{Turing.FinTM.bufferTape, Turing.FinTM.bufferTape_nat}
\difficulty{3}
Let $w$ be a Boolean word and let $i$ be a natural number with $i<|w|$. The tape storing
$w$ starting at cell zero contains the symbol $w_i$ at cell $i$.
\end{theorem}

\begin{theorem}[A Sweep Across a Boolean Buffer]
\lean{Complexity.Meyer.dec_sweep}
\label{Complexity.Meyer.dec_sweep}
\leanok
\uses{Complexity.Meyer.bufferTape_at, Complexity.Meyer.dec_zero, Complexity.Meyer.lcfg, Complexity.Meyer.tabTM, Complexity.TimeHierarchy.OutReg, Complexity.TimeHierarchy.bufferTape_update_mid, Turing.Cfg, Turing.FinTM.bufferTape, Turing.MultiTapeTM.runFrom, Turing.MultiTapeTM.runFrom_succ_eq_step}
\difficulty{5}
Let $M$ be a finite Boolean Turing machine, let $y,x$ be Boolean words, let $bits$ be a
selector-bit assignment for $M$, let $P$ be a valid input-head position for $y$, let $c$
be a configuration of $M$ on input $x$, and let $\tau_O$ be an integer-indexed Boolean
buffer with head position $z_O$. For all $i,j\in\mathbb{N}$ and all Boolean words
$rest$, starting the simulator in its decode state with $C$ buffer containing $i$ true
bits followed by $j$ false bits and then $rest$, with its head at position $i$, after
$j$ steps the $C$ buffer contains $i+j$ true bits followed by $rest$ and its head is at
position $i+j$; all the other listed configuration data are unchanged.
\end{theorem}

\begin{theorem}[Return-State Sweep]
\lean{Complexity.Meyer.ret_sweep}
\label{Complexity.Meyer.ret_sweep}
\leanok
\uses{Complexity.Meyer.bufferTape_at, Complexity.Meyer.lcfg, Complexity.Meyer.ret_none, Complexity.Meyer.ret_some, Complexity.Meyer.tabTM, Complexity.TimeHierarchy.OutReg, Turing.Cfg, Turing.FinTM.bufferTape, Turing.MultiTapeTM.runFrom, Turing.MultiTapeTM.runFrom_succ_eq_step, Turing.MultiTapeTM.runFrom_zero}
\difficulty{5}
Let $M$ be a finite Turing machine over the binary alphabet, let $x$ and $y$ be bit
strings, let $c$ be a configuration of $M$ on input $x$, and fix a selector bit vector,
an input-head position on $y$, a return state of $M$, an output-register value, a
Boolean boundary tag, an auxiliary output tape with head position $z_O$, and a bit
string $L$. For every natural number $j$, starting in the corresponding return-state
configuration with the first auxiliary tape storing $j$ consecutive ones followed by $L$
and its head at position $j-1$, the tabulating machine reaches after $j+1$ steps the
corresponding decision-state configuration with that tape's head at position $0$; the
tape contents and all other specified data are unchanged.
\end{theorem}

\begin{theorem}[Bounded simulation of a machine run]
\lean{Complexity.Meyer.loop_run}
\label{Complexity.Meyer.loop_run}
\leanok
\uses{Complexity.Meyer.bufferTape_at, Complexity.Meyer.dec_fire, Complexity.Meyer.dec_sweep, Complexity.Meyer.dec_under, Complexity.Meyer.lcfg, Complexity.Meyer.ret_sweep, Complexity.Meyer.tabTM, Complexity.TimeHierarchy.OutReg.ofList, Complexity.TimeHierarchy.bufferTape_update_mid, Complexity.TimeHierarchy.ctrVal, Complexity.TimeHierarchy.ctrVal_borrow, Complexity.TimeHierarchy.ctrVal_pos_decomp, Complexity.TimeHierarchy.ctrVal_zero_eq, Turing.Cfg, Turing.FinTM.VirtualTag, Turing.FinTM.bufferTape, Turing.MultiTapeTM.runFrom, Turing.MultiTapeTM.runFrom_add, Turing.MultiTapeTM.runFrom_of_halt, Turing.MultiTapeTM.runFrom_succ_eq_step, Turing.MultiTapeTM.runFrom_succ_eq_step', Turing.MultiTapeTM.step}
\difficulty{6}
Let $M$ be a deterministic Turing machine with Boolean symbols, and let $x,y,w$ be
finite Boolean words. Let $P$ be an input-head position for $y$, let $bits$ assign a
Boolean value to each of the $nbits(M)$ positions, and let
$\tau_O:\mathbb{Z}\to\mathrm{Option}(\mathrm{Bool})$ and $z_O\in\mathbb{Z}$. For every
natural number $v$, configuration $c$ of $M$ on input $x$, state $q$ of $M$, and Boolean
tag $b$ such that $ctrVal(w)=v$, the state of $c$ is $q$, and $b$ satisfies the
boundary-tag condition for the input-head position of $c$, there are a natural number
$s\le (v+1)(2|w|+3)$, a tape content function
$\tau_C:\mathbb{Z}\to\mathrm{Option}(\mathrm{Bool})$, an integer head position $z_C$,
and a Boolean tag $b'$ such that $b'$ satisfies the boundary-tag condition for the
input-head position of the configuration reached by $M$ from $c$ after $v$ steps. After
$s$ steps, the tabulated machine started from the specified encoded configuration in its
decode state reaches the specified encoded configuration in its write-decode state,
carrying the machine's state and output register after those $v$ steps, its resulting
configuration, and the tag $b'$, while using $\tau_C,z_C$ for the auxiliary tape and
head and retaining $\tau_O,z_O$.
\end{theorem}
